% cosmology — полный том «Космология и наблюдаемый макромир» (book/SERIES.md)
\part{Космология каузального пузыря над кипящим океаном}

\chapter{Два времени: вечно кипящий океан микроуровня \texorpdfstring{$\Mlayer$}{M} и эволюция каузального пузыря на макроуровне \texorpdfstring{$\Tlayer$}{T}}
\label{ch:cosmo-two-times}

В континуальной космологии понятие «Вселенная» объединяет в один объект и саму физическую среду (вакуум), и наблюдаемую совокупность галактик и реликтового излучения.
Такое отождествление неизбежно порождает парадокс начальной сингулярности ($t = 0$, $\rho \to \infty$, рождение самого пространства-времени из небытия) и парадокс космологической постоянной (расхождение между планковской плотностью энергии вакуума $u_P \approx 4{,}63\times 10^{113}\ \NotationUnitJoulePerCubicMetre$ и наблюдаемой кривизной горизонта на $122$ порядка).

В дискретной теории решёточной жидкости на гранецентрированной кубической решётке~\latGCK\ эти два объекта строго разделены по уровням описания:
\begin{enumerate}[leftmargin=*]
\item \textbf{Первичный носитель (кипящий океан микроуровня~$\Mlayer$):}
  фиксированная трёхмерная решётка~$\carrierLattice$ с шагом $\caStepSpace = \ell_P$, тактом $\htick = t_P/\sqrt{2}$ и обратимым унитарным шагом эволюции~$g$ на кольце спиноров $\mathbb{Z}_{\phaseRingCount}[i]^2$ ($\phaseRingCount = 2^{\lfloor\caBitBudget\rfloor} = 512$).
  В силу закона сохранения полной нормы ($A_3$, $\sum_x |z(x)|^2 = \mathrm{const}$), обратимости шага ($A_{13}$, существование~$g^{-1}$) и недостижимости абсолютного нуля ($A_5$, $|z(x)| \ge z_{\min} = 2^{-\mathrm{frac\_bits}} = 1/64$) первичный океан \emph{не имеет ни начала во времени, ни тепловой смерти}.
  Решётка не расширяется и не сжимается; её узлы не разлетаются в пустоте, а шаг~$\ell_P$ остаётся неизменным на всех этапах космологической эволюции.
\item \textbf{Наблюдаемый каузальный пузырь (макроуровень~$\Tlayer$):}
  локализованная область фазового возбуждения внутри вечно кипящего океана, возникшая $N_{\mathrm{today}} \approx 1{,}141\times 10^{61}$ тактов~$\htick$ назад ($t_{\mathrm{age}} = N_{\mathrm{today}}\,\htick \approx 13{,}787\ \mathrm{млрд\ лет}$) в результате детерминированного фазового сдвига синхронного «льда».
  Все наблюдаемые космологические явления --- расширение каузального горизонта $R_H(\tau) = c\,\tau\,\htick = \caKappaGeom\,\tau\,\ell_P$, нуклеосинтез, рекомбинация на такте $N_{\mathrm{CMB}} \approx 3{,}146\times 10^{56}$ ($t_{\mathrm{recomb}} \approx 380\ \mathrm{тыс.\ лет}$) и формирование крупномасштабной структуры --- представляют собой распространение и релаксацию волн и вихревых солитонов по неподвижной решётке~$\carrierLattice$.
\end{enumerate}

\begin{theorem}[Мост времени $t_{\mathrm{SI}}(\tau) = \tau\,\htick$ и отсутствие сингулярности]
\label{thm:cosmo-time-bridge}
\begin{enumerate}[leftmargin=*]
\item \textbf{Счётчик тактов каузального пузыря:}
  при выборе локального отсчёта $\tau = 0$ в момент старта расширения нашего каузального пузыря (переход Фаза~IV $\to$ Фаза~I) физическое время любого космологического события в секундах СИ связано с целым числом тактов автомата $\tau \in \mathbb{Z}_{\ge 0}$ точным соотношением без свободных параметров:
  \begin{equation}
  \label{eq:cosmo-exact-tick-time}
    t_{\mathrm{SI}}(\tau) = \tau\,\htick = \tau\,\caKappaGeom\,t_P,
    \qquad
    \htick = \frac{\caStepSpace}{\caSignalSpeed} = \frac{t_P}{\sqrt{2}} \approx 3{,}812\times 10^{-44}\ \mathrm{с}.
  \end{equation}
  Современному возрасту пузыря $t_{\mathrm{age}} \approx 13{,}787\ \mathrm{млрд\ лет}$ и эпохе последней поверхности рассеяния реликтового излучения $t_{\mathrm{recomb}} \approx 380\ \mathrm{тыс.\ лет}$ отвечают целые числа тактов:
  \[
    N_{\mathrm{today}} = \left\lfloor\frac{t_{\mathrm{age}}}{\htick}\right\rfloor \approx 1{,}141\times 10^{61},
    \qquad
    N_{\mathrm{CMB}} = \left\lfloor\frac{t_{\mathrm{recomb}}}{\htick}\right\rfloor \approx 3{,}146\times 10^{56}.
  \]
\item \textbf{Планковский монолит вместо начальной сингулярности ($\rho \le \rho_P$):}
  поскольку в каждой ячейке~$V_P$ занятость веществом ограничена булевым потолком $b(x) \in \{0,1\}$, а плотность энергии спинора ограничена планковской упругостью $K_P = u_P \approx 4{,}63\times 10^{113}\ \mathrm{Па}$ ($A_7$), при $\tau \to 0$ плотность энергии и кривизна Редже нигде не обращаются в бесконечность ($\rho_{\mathrm{matter}} \le \rho_P \approx 5{,}155\times 10^{96}\ \NotationUnitKgPerMetreCubed$).
\end{enumerate}
\end{theorem}

\chapter{Девятибитный начальный вектор Вселенной (\texorpdfstring{$\mathrm{IV}\in\mathbb{Z}_{512}$}{IV in Z512}), механизм Киббла--Зурека и барионная асимметрия}
\label{ch:cosmo-9bit-iv}

Поскольку микроскопический шаг автомата~$g$ строго обратим ($A_{13}$), вся история нашего каузального пузыря однозначно кодируется его начальным состоянием в момент нуклеации $\tau = 0$.
В силу неинъективности макроскопического усреднения~$\mathcal{B}$ ($A_6$) обратить полную микроскопическую волновую функцию $\Psi(\tau)$ по одной сегодняшней карте реликтового излучения невозможно (эффект информационной завесы).
Однако конечность битового бюджета планковской ячейки Бекештейна позволяет поставить и решить обратную задачу для \emph{начального вектора зародыша пузыря} ($\mathrm{IV}$).

\section{Четыре фазы информационного цикла океана}
\label{sec:cosmo-four-phases}

На тороидальной или неограниченной решётке~\latGCK\ с сохранением нормы~$A_3$ и ненулевым полом кипения~$A_5$ эволюция локальной области проходит четыре закономерные фазы без изменения шага решётки:
\begin{enumerate}[leftmargin=*]
\item \textbf{Фаза~I (Расширение и рождение вихрей --- наша эпоха):}
  локальный фазовый сдвиг на фоне синхронного океана запускает каскад нелинейного затвора насыщающей фазы~$A_7$. Фронт возбуждения бежит по каузальной звезде~$\geomNeighborTwelve$ со скоростью света $c = \caKappaGeom\caSignalSpeed$, а на стыках фазовых доменов рождаются топологически защищённые вихревые солитоны ($b=1$, $n\in\mathbb{Z}$), формирующие барионы, лептоны, ядра и галактики.
\item \textbf{Фаза~II (Макроскопическое затухание на уровне~$\Tlayer$):}
  по мере рассеяния и аннигиляции встречных вихрей ($n=+1$ и $n=-1 \to 2\gamma$) макроскопический контраст плотности падает, а энергия макро-движения перекачивается через решёточную вязкость $\nu_{\mathrm{CA}} = \kappa_{\mathrm{link}}\caSignalSpeed\ell_P$ в мелкомасштабное фазовое кипение. При этом на микроуровне~$\Mlayer$ полная норма $\sum_x |z(x)|^2$ сохраняется тождественно (тепловая смерть существует только для грубого прибора на~$\Tlayer$).
\item \textbf{Фаза~III (Синхронный «лёд»):}
  область релаксирует к однородному синхронному вращению с минимальным квантом амплитуды $z_{\min} = 2^{-\mathrm{frac\_bits}}$ и нулевым относительным градиентом классов фазы между соседями ($\Delta\mathrm{class}_{N} = 0$). Для макроскопического прибора такое состояние неотличимо от абсолютной пустоты, тогда как на~$\Mlayer$ в каждой ячейке продолжается тактовое вращение.
\item \textbf{Фаза~IV (Нуклеация нового пузыря):}
  накопление детерминированной голономии Гейзенберга ($\Delta\varphi_{\mathrm{disc}} = 41$ такт кольца) на границе синхронного домена создаёт локальный трёхгранный фазовый сдвиг, запускающий новый цикл Фазы~I.
\end{enumerate}

\section{Теорема о 9-битном начальном векторе зародыша (\texorpdfstring{$\mathrm{IV}=263$}{IV=263}) и барионной асимметрии \texorpdfstring{$\eta_B$}{eta\_B}}
\label{sec:cosmo-9bit-theorem}

Почему плоская доменная стенка или сферически-симметричный сдвиг общей фазы не рождают вещества ($b=0$), а наш каузальный пузырь заполнен барионной материей с преобладанием положительного топологического заряда ($n=+1$) и точным отношением числа барионов к числу фотонов $\eta_B = n_B/n_\gamma \approx 6{,}12\times 10^{-10}$?

\begin{theorem}[Девятибитный сид Вселенной $\mathrm{IV}\in\mathbb{Z}_{512}$ и барионная асимметрия]
\label{thm:cosmo-9bit-iv-baryon}
\begin{enumerate}[leftmargin=*]
\item \textbf{Точное совпадение размерности начального вектора с битовым бюджетом ячейки ($\dim_2(\mathrm{IV}) = \lfloor\caBitBudget\rfloor = 9\text{ бит}$):}
  вокруг центрального узла нуклеации $(c_z, c_y, c_x)$ три взаимно ортогональные координатные плоскости ($x\ge c_x$, $y\ge c_y$, $z\ge c_z$), несущие три цветовых цикла барионного треугольника~$K_3$, делят трёхмерное пространство ровно на
  \[
    2^d = 2^3 = 8 = \vevHierarchySteps\text{ пространственных октантов.}
  \]
  Бинарная маска начального фазового сдвига Гейзенберга $\mathrm{mask}_8 \in \{0, 1, \ldots, 255\}$ на этих $\vevHierarchySteps = 8$ октантах в сочетании с одним старшим битом хиральности спинора $b_{\mathrm{chir}} \in \{0, 1\}$ (знак относительного фазового сдвига между компонентами $z_1$ и~$z_2$) образует \textbf{9-битный начальный вектор зародыша Вселенной}:
  \begin{equation}
  \label{eq:cosmo-9bit-iv-def}
  \begin{aligned}
    \mathrm{IV} &= (b_{\mathrm{chir}} \ll 8) \;\big|\; \mathrm{mask}_8 \;\in\; \bigl\{0, 1, \ldots, \phaseRingCount - 1\bigr\},\\
    1_{\mathrm{chir}} + 2^d_{\mathrm{octants}} &= 1 + \vevHierarchySteps = \lfloor\caBitBudget\rfloor = 9\text{ бит}\quad(\phaseRingCount = 2^9 = 512).
  \end{aligned}
  \end{equation}
\item \textbf{Трёхгранный узел Киббла--Зурека и канонический сид нашего пузыря ($\mathrm{IV}_{\mathrm{bubble}} = 263 = \texttt{0b1\_00000111}$):}
  при однородной маске ($\mathrm{mask}_8 = 0$ или $255$) либо плоской одномерной стенке контурный интеграл фазы вокруг любой замкнутой петли равен нулю ($\oint\mathrm{d}\arg z = 0$), и вихри материи не рождаются ($b = 0$).
  Нетривиальная топологическая голономия возникает тогда и только тогда, когда вокруг общей вершины сходятся три ортогональных октанта барионного угла~$K_3$ ($\mathrm{mask}_8 = \texttt{0b00000111} = 7 = d_{\mathrm{soft}}$) при положительном бите хиральности $b_{\mathrm{chir}} = 1$:
  \begin{equation}
  \label{eq:cosmo-canonical-seed-263}
    \mathrm{IV}_{\mathrm{bubble}} = 2^{\vevHierarchySteps} + d_{\mathrm{soft}} = 256 + 7 = 263 = \texttt{0b1\_00000111}.
  \end{equation}
  Прямое численное интегрирование обратимого шага~$g$ на трёхмерной решётке~\latGCKshort\ с начальным вектором $\mathrm{IV} = 263$ из состояния без посаженных вихрей ($b(\tau=0) \equiv 0$) детерминированно рождает устойчивые вихревые солитоны материи ($b = 1$) с чистым положительным топологическим зарядом $\sum n > 0$, тогда как инверсия старшего бита ($\mathrm{IV} = 7$, $b_{\mathrm{chir}} = 0$) зеркально меняет знак рождающегося мира на антиматерию ($\sum n < 0$).
\item \textbf{Количественный вывод барион-фотонного отношения $\eta_B = n_B / n_\gamma$:}
  при прохождении фронта нуклеации через $\vevHierarchySteps = 8$ ступеней электрослабой башни доля выживших барионных узлов~$K_3$ по отношению к числу родившихся безмассовых квантов излучения ($n=0$, фотонов реликтовой ванны) определяется четвёртой степенью константы сцепления~$\alpha_{\mathrm{fs}}^4$, умноженной на CP-нечётную мнимую часть вершины треугольника унитарности $\bar\eta = \lfloor\caBitBudget\rfloor / (2N_\varphi) = 9/26$ и долю внешних рёберных каналов $d_{\mathrm{out}}(E)/\vevHierarchySteps = 5/8$:
  \begin{equation}
  \label{eq:cosmo-baryon-to-photon}
    \eta_B := \frac{n_B}{n_\gamma}
    = \bar\eta\,\frac{d_{\mathrm{out}}(E)}{\vevHierarchySteps}\,\alpha_{\mathrm{fs}}^4
    = \frac{45}{208}\,\alpha_{\mathrm{fs}}^4
    \approx 6{,}135\times 10^{-10}
  \end{equation}
  (наблюдение Planck~2018: $(6{,}12 \pm 0{,}04)\times 10^{-10}$, откл.~$0{,}24\%$).
\end{enumerate}
\end{theorem}

\chapter{Структура космологического горизонта: \texorpdfstring{$\Lambda_{\mathrm{bubble}}$}{Lambda}, отношение тёмного ореола к барионам \texorpdfstring{$\Omega_{\mathrm{DM}}/\Omega_b = 16/3$}{16/3} и тепловой спектр пузыря}
\label{ch:cosmo-horizon-dm-cmb}

\section{Решение проблемы космологической постоянной: \texorpdfstring{$\Lambda_{\Mlayer} \equiv 0$}{Lambda\_M = 0} и кривизна горизонта \texorpdfstring{$\Lambda_{\mathrm{bubble}}$}{Lambda\_bubble}}
\label{sec:cosmo-lambda-solution}

В главе~\bref{ch:si-sm} тома «Построение» (\bcref{rem:channel-gr-open-in-book}) установлено, что макроскопическая гравитация представляет собой упругую деформацию Редже $\varepsilon_e = (\rho_e - \rho_{\mathrm{vac}})/\rho_{\mathrm{vac}}$ длин связей относительно однородного кипящего фона~$A_5$.
Отсюда вытекают два фундаментальных космологических следствия:
\begin{enumerate}[leftmargin=*]
\item \textbf{Тождественное зануление вклада вакуумного океана ($\Lambda_{\Mlayer} \equiv 0$):}
  в однородном кипящем океане $\rho_e \equiv \rho_{\mathrm{vac}} \implies \varepsilon_e \equiv 0 \implies \delta_e \equiv 0 \implies G_{\mu\nu}[\text{однородный океан}] \equiv 0$.
  Фоновая плотность энергии кипения $\rho_{E,\mathrm{vac}} = 2z_{\min}^2 u_P$ является точкой отсчёта ненапряжённой решётки и не искривляет макроскопическую метрику, что полностью устраняет ошибку в $10^{122}$ раз континуальной квантовой теории поля.
\item \textbf{Горизонтная кривизна конечного каузального пузыря $\Lambda_{\mathrm{bubble}}$ и сферический фактор проектора~$\mathcal{B}$:}
  на такте $N_{\mathrm{today}} \approx 1{,}141\times 10^{61}$ от момента нуклеации свет со скоростью $c = \caKappaGeom\caSignalSpeed$ прошёл конечный каузальный радиус $R_H = \caKappaGeom N_{\mathrm{today}}\ell_P$.
  Неизотропный кубический потолок по $|\square| = d/\caKappaGeom^2 = 6$ квадратным граням даёт $\Lambda_{\mathrm{raw}}\,\ell_P^2 = |\square|/N_{\mathrm{today}}^2 \approx 4{,}61\times 10^{-122}$.
  При переходе через макропроектор~$\mathcal{B}$ от гранёного однотактового кубооктаэдра (полная площадь граней которого складывается из $|\square|=6$ квадратов площади~$a^2$ и $|\triangle|=8$ правильных треугольников площади $\frac{\sqrt{3}}{4}a^2$, то есть $A_{\mathrm{cuboct}} = (6 + 2\sqrt{3})a^2 = 2(d+\sqrt{d})a^2$) к изотропной сфере горизонта радиуса $R_{\mathrm{in}} = \caKappaGeom a = a/\sqrt{2}$ (площадь $A_{\mathrm{sphere}} = 4\pi R_{\mathrm{in}}^2 = 2\pi a^2$) действует точный геометрический коэффициент сферической проекции:
  \[
    \eta_{\mathrm{surf}} = \frac{A_{\mathrm{sphere}}}{A_{\mathrm{cuboct}}} = \frac{2\pi a^2}{2(d+\sqrt{d})a^2} = \frac{\pi}{d + \sqrt{d}} = \frac{\pi}{3 + \sqrt{3}} \approx 0{,}6639,
  \]
  откуда наблюдаемая горизонтная кривизна пузыря в планковских единицах равна:
  \begin{equation}
  \label{eq:cosmo-lambda-bubble}
    \Lambda_{\mathrm{bubble}}\,\ell_P^2
    = \eta_{\mathrm{surf}}\,\frac{|\square|}{N_{\mathrm{today}}^2}
    = \frac{3 - \sqrt{3}}{2}\,\frac{2\pi}{N_{\mathrm{today}}^2}
    \approx 3{,}06\times 10^{-122}
  \end{equation}
  (наблюдаемое значение Planck~2018: $\Lambda_{\mathrm{obs}}\,\ell_P^2 \approx (2{,}89 \pm 0{,}09)\times 10^{-122}$, $E_n \approx 0{,}94 < 2$).
\end{enumerate}

\section{Запрет частиц «тёмной материи» (\texorpdfstring{$D_{\mathrm{dark}} \equiv 0$}{D\_dark = 0}) и вакуумно-голономный ореол \texorpdfstring{$\Omega_{\mathrm{halo}}/\Omega_b = 16/3$}{16/3}}
\label{sec:cosmo-dark-matter-ratio}

В астрофизических наблюдениях (плоские кривые вращения галактик $v(r)\to\mathrm{const}$, гравитационное линзирование скоплений и акустические пики реликтового излучения Planck~2018) полная гравитационная деформация вокруг барионных скоплений систематически превышает ньютоновский вклад локализованных светящихся барионов в $\Omega_{\mathrm{DM}}/\Omega_b = \Omega_c h^2 / \Omega_b h^2 = 0{,}1200 / 0{,}02237 \approx 5{,}364 \pm 0{,}06$ раза.
В континуальной космологии~$\Lambda\mathrm{CDM}$ этот избыток приписывается гипотетическим калибровочно-стерильным частицам «тёмной материи» (WIMP, аксионам, скрытому сектору).
На дискретном носителе, однозначно определяемом единственным корнем уравнения Каталана $1+2^d=d^2 \iff d=3$ (\texttt{catalan\_\allowbreak lock\_\allowbreak iff\_\allowbreak d3}), существование подобных частиц математически невозможно, а сам эффект гало и отношение $16/3$ выводятся как строгое целочисленное тождество:

\begin{theorem}[Запрет частиц «тёмной материи» $D_{\mathrm{dark}}\equiv 0$, вакуумно-голономная природа гало на полу $A_5$ и отношение $\Omega_{\mathrm{halo}}/\Omega_b = 16/3$]
\label{thm:cosmo-no-dark-matter-and-halo-ratio}
\label{prop:cosmo-dm-baryon-ratio}
Пусть носитель фиксирован единственным решением уравнения Каталана $d=3$ ($9$-битное кольцо $\mathbb{Z}_{512}$, кубооктаэдр $O_2$, пол вакуума $\mu_{\mathrm{vac}}=1$, потолок $K_P=64$, $\varepsilon=8192$). Тогда выполняются четыре точных свойства (доказанных и машинно проверенных в Lean~4 как теорема \texttt{no\_\allowbreak dark\_\allowbreak matter\_\allowbreak and\_\allowbreak vacuum\_\allowbreak holonomy\_\allowbreak halo} в модуле \texttt{UnifiedOperators.lean}):
\begin{enumerate}[leftmargin=*]
\item \textbf{Тождество полноты $4$-канального носителя и отсутствие «тёмного остатка» ($D_{\mathrm{dark}}\equiv 0$):}
  в силу тождества чётной голономии $C(d\varphi_{xy}) + |A_{xy}| = 128$ (\texttt{even\_\allowbreak holonomy\_\allowbreak bohm\_\allowbreak deficit\_\allowbreak eq\_\allowbreak abs\_\allowbreak connection}, где $A_{xy} = \mathrm{wrap}_{512}(\varphi_y - \varphi_x)$) для любых целочисленных амплитуд $\mu_x,\mu_y\in\mathbb{Z}$ и любых фаз гравитационный знаменатель преломления фазовых часов $D(x) = \zeta_{\mathrm{real}}(x) + 128\,\mu_x^2 + 8192$ раскладывается \emph{тождественно и без остатка} на чисто барионный член и наблюдаемый в дисперсионном канале дефицит Бома $\delta K_{xy} = \mu_x\mu_y|A_{xy}|$:
  \begin{equation}
  \label{eq:cosmo-zero-dark-residual}
    D(x) \;-\; \Bigl(\underbrace{8192 + 128\,\mu_x^2 + 128\,\mu_x\sum_{y\in x+O_2}\mu_y}_{D_{\mathrm{baryon}}(\mu)} \;-\; \underbrace{\sum_{y\in x+O_2}\mu_x\mu_y\,|A_{xy}|}_{\delta K(x)}\Bigr) \;\equiv\; 0.
  \end{equation}
  В $9$-битном регистре Каталана $\mathbb{Z}_{512}$ не существует ни одного калибровочно-стерильного («тёмного») бита: остаток $D_{\mathrm{dark}}(x) \equiv 0$ на всех конфигурациях решётки.
\item \textbf{Необходимость единичного пола вакуума~$A_5$ ($\mu_{\mathrm{vac}}=1$) и бесчастичная природа галактических гало:}
  в гипотетическом классическом пустом вакууме ($\mu_x=\mu_y=0$) знаменатель $D(0,0,A) \equiv 8192$ вообще не зависит от фазовой связности~$A_{xy}$ ($\delta K_{xy}\equiv 0$).
  Напротив, на доказанном в теореме \texttt{a5\_\allowbreak vacuum\_\allowbreak floor\_\allowbreak from\_\allowbreak ring\_\allowbreak visibility\_\allowbreak and\_\allowbreak iv263} единичном полу вакуума ($\mu_x=\mu_y=\mu_{\mathrm{vac}}=1$, отсутствие барионного вещества) знаменатель импеданса равен:
  \begin{equation}
  \label{eq:cosmo-vacuum-halo-impedance}
    D_{\mathrm{vac}}(x) \;=\; 9856 \;-\; \sum_{y\in x+O_2} |A_{xy}|,
    \qquad
    D_{\mathrm{vac}}(x) \ne 9856 \;\Longleftrightarrow\; \sum_{y\in x+O_2} |A_{xy}| > 0.
  \end{equation}
  Таким образом, любое гравитационное линзирование и центростремительное преломление $\Delta_t A_e = (d_0\Phi)_e$ в безбарионном ореоле галактик ($\mu=1$) порождается не «облаком тёмных частиц», а чисто фазовым натяжением $|A_{xy}| > 0$ единичного вакуумного пола~$A_5$ вокруг барионного вихревого ядра (на пороге Гейзенберга $|A_{xy}|=41$ тик дающим $D_{\mathrm{vac}}=9364$ и выход кривых вращения на плато $v^2(r)=\mathrm{const}$, а при столкновении скоплений типа Bullet Cluster проходящим сквозь заторможенный потолком $A_7/A_{16}$ горячий газ $\mu>1$ вместе с вихревыми ядрами галактик).
\item \textbf{Точное отношение ёмкости кубооктаэдрального ореола к $3\mathrm{D}$-ядру Вороного $\Omega_{\mathrm{halo}}/\Omega_b = 16/3$ из одного корня $d=3$:}
  целочисленная голономная ёмкость $12$-лучевого кубооктаэдрального ореола деформации равна $W_{\mathrm{halo}} = 2\,|\triangle(3)|\,Q_\square(3) = 2\cdot 2^3\cdot 128 = 2048$ (чему в геометрической мере при ребре $a=\ell_P$ отвечает объём кубооктаэдра $V_{\mathrm{cuboct}} = \frac{8\sqrt{2}}{3}a^3 = \frac{2\vevHierarchySteps}{d}\caKappaGeom a^3$), тогда как ёмкость локализованного $3\mathrm{D}$-ядра Вороного равна $W_{\mathrm{core}} = d\,Q_\square(3) = 3\cdot 128 = 384$ (объём ячейки Вороного--Вигнера--Зейтца $dV_{\mathrm{Voronoi}} = a^3/\sqrt{2} = \caKappaGeom a^3$). Их отношение есть точная несократимая константа корня $d=3$:
  \begin{equation}
  \label{eq:cosmo-dm-baryon-16-over-3}
    \frac{\Omega_{\mathrm{DM}}}{\Omega_b}
    = \frac{W_{\mathrm{halo}}}{W_{\mathrm{core}}}
    = \frac{2048}{384}
    = \frac{V_{\mathrm{cuboct}}}{dV_{\mathrm{Voronoi}}}
    = \frac{2\,\vevHierarchySteps}{d}
    = \frac{16}{3} \approx 5{,}3333
  \end{equation}
  (Planck~2018: $5{,}364 \pm 0{,}06$, отклонение~$0{,}58\%$), а отношение полной гравитационной деформации (ядро $+$ ореол) к барионной равно $(W_{\mathrm{halo}}+W_{\mathrm{core}})/W_{\mathrm{core}} = 2432/384 = (2\vevHierarchySteps + d)/d = 19/3 \approx 6{,}333$ (Planck~2018: $6{,}364$, откл.~$0{,}48\%$).
\item \textbf{Точный нулевой прогноз для детекторов частиц тёмной материи ($N_{\mathrm{DM\text{-}particles}} \equiv 0$):}
  поскольку в ореоле галактик $\mu = \mu_{\mathrm{vac}} = 1$ и $D_{\mathrm{dark}} \equiv 0$, в природе не существует корпускулярных носителей тёмной материи ($\mu > 1$, $\delta K_{xy} = 0$), способных вызывать ядерную или электронную отдачу в подземных, коллайдерных и аксионных детекторах; прямое детектирование гало возможно исключительно через фазово-голономный канал Бома--Ааронова $\delta K_{\mathrm{vac}} = |A_{xy}| = 9856 - D_{\mathrm{vac}}$.
\end{enumerate}
\end{theorem}

\section{Две тепловые ванны и замкнутый вывод теплового кванта реликтового излучения в естественной шкале (без \texorpdfstring{$k_{\mathrm{B}}$}{k\_B})}
\label{sec:cosmo-two-baths}

Разделение уровней~$\Mlayer$ и~$\Tlayer$ устраняет смешение микроскопического кипения носителя с макроскопическим фотонным фоном пузыря и позволяет вычислить тепловой квант реликтового излучения непосредственно в безразмерной энергетической шкале лестницы (где постоянная Больцмана~$k_{\mathrm{B}}$ выступает исключительно коэффициентом перевода энергетических единиц в градусы Кельвина на приборном слое~$\Tlayer$):
\begin{itemize}[leftmargin=*]
\item \textbf{Микроскопический тепловой квант кипящего пола~$A_5$ в шкале носителя ($\Theta_{\Mlayer,\mathrm{vac}}/E_0$ без~$\pi$):}
  при минимальном кванте амплитуды $z_{\min} = a_Q = 2^{-\mathrm{frac\_bits}}$ на каждой из $d-1$ компонент спинора фоновая плотность энергии океана равна $\rho_{E,\mathrm{vac}} = (d-1)a_Q^2 u_P = \frac{1}{d_{\mathrm{in}}(\square)\,\phaseRingCount}\,u_P$. Поскольку тактовый квант носителя равен $E_0 = s_0/\htick = \caKappaGeom E_P$ ($E_P^{d_{\mathrm{in}}(\square)} = d_{\mathrm{in}}(\square)\,E_0^{d_{\mathrm{in}}(\square)}$), в естественной шкале~$E_0$ множитель $d_{\mathrm{in}}(\square)$ сокращается нацело, давая чисто алгебраический инвариант кольца без~$\pi$:
  \[
    \frac{\rho_{E,\mathrm{vac}}}{E_0^{d_{\mathrm{in}}(\square)}/(\planckHbar c)^d} = \frac{1}{\phaseRingCount}
    \quad\Longrightarrow\quad
    \frac{\Theta_{\Mlayer,\mathrm{vac}}}{E_0} = \phaseRingCount^{-1/d_{\mathrm{in}}(\square)} \approx 0{,}2102
  \]
  (круговой множитель~$\pi^2$ возникает только при переходе к континуальному интегралу Стефана--Больцмана на~$\Tlayer$, где $T_{\Mlayer,\mathrm{bath}} \sim 10^{31}\ \mathrm{К}$). Эта энергия удерживается на собственной частоте носителя $\nu_0 = 1/\htick$ на границе зоны Бриллюэна.
\item \textbf{Естественный тепловой квант рекомбинации ($\Theta_{\mathrm{rec}}$) и современного реликтового фона ($\Theta_{\mathrm{CMB}}$):}
  длинноволновый бозевский газ свободных фотонов ($n=0$, $\omega = c|\mathbf{k}| \ll \omega_{\max}$) отделяется от барионной плазмы на такте рекомбинации~$N_{\mathrm{CMB}}$ и остывает при расширении каузального фронта пузыря до современного такта~$N_{\mathrm{today}}$.
\end{itemize}

\begin{theorem}[Алгебраический вывод тепловых квантов рекомбинации $\Theta_{\mathrm{rec}}$ и реликтового фона $\Theta_{\mathrm{CMB}}$ в шкале носителя без $\pi$, без $\ln$ и без $k_{\mathrm{B}}$]
\label{thm:cosmo-cmb-natural-scale}
В естественной энергетической шкале носителя $T_{P,\Mlayer} := E_0 = s_0/\htick$ (без переводного коэффициента~$k_{\mathrm{B}}$, без~$\pi$ и без голых численных констант) тепловой квант отцепления фотонов при рекомбинации водорода~$\Theta_{\mathrm{rec}}$ и современный тепловой квант реликтового излучения~$\Theta_{\mathrm{CMB}}$ однозначно определяются комбинаторикой кубооктаэдра и постоянной тонкой структуры~$\alpha_{\mathrm{fs}}$:
\begin{enumerate}[leftmargin=*]
\item \textbf{Алгебраический порог отцепления Лайман-$\alpha$ и квант рекомбинации $\Theta_{\mathrm{rec}}$:}
  энергия главного дипольного перехода водорода равна $E_{\mathrm{Ly}\alpha} = \frac{d}{d_{\mathrm{in}}(\square)}\,E_{\mathrm{Ry}}$, где $E_{\mathrm{Ry}} = \frac{\alpha_{\mathrm{fs}}^{d-1}}{d-1}\,m_e c^2$.
  На дискретном носителе отцепление фотонной бани от нейтрального водорода наступает тогда, когда квант Лайман-$\alpha$ распределяется по $d_{\mathrm{out}}(E)\,\vevHierarchySteps$ внешним иерархическим рёберным каналам кубооктаэдра ($E=(2,5,5)$, $d_{\mathrm{out}}(E)=5$):
  \begin{equation}
  \label{eq:cosmo-theta-rec-natural}
    \frac{\Theta_{\mathrm{rec}}}{E_{\mathrm{Ly}\alpha}}
    = \frac{1}{d_{\mathrm{out}}(E)\,\vevHierarchySteps}
    \quad\Longleftrightarrow\quad
    \frac{\Theta_{\mathrm{rec}}}{E_{\mathrm{Ry}}}
    = \frac{d}{d_{\mathrm{in}}(\square)\,d_{\mathrm{out}}(E)\,\vevHierarchySteps}
    \approx 0{,}01875
  \end{equation}
  (в континуальном описании Саха на~$\Tlayer$ тому же отношению отвечает логарифм барион-фотонного отношения $\bigl[\frac{d_{\mathrm{out}}(E)}{d-1}\ln(1/\eta_B)\bigr]^{-1} \approx 0{,}018857$ при $\eta_B^{-1} = \frac{(d-1)N_\varphi\vevHierarchySteps}{\lfloor\caBitBudget\rfloor d_{\mathrm{out}}(E)}\,\alpha_{\mathrm{fs}}^{-d_{\mathrm{in}}(\square)}$, сходящийся с алгебраическим отношением носителя~\eqref{eq:cosmo-theta-rec-natural} с точностью~$0{,}57\%$).
  При переводе в шкалу Кельвина на приборном слое~$\Tlayer$ через $T_{\mathrm{Ry}} = E_{\mathrm{Ry}}/k_{\mathrm{B}}$ алгебраическая формула~\eqref{eq:cosmo-theta-rec-natural} даёт $T_{\mathrm{rec}}^{\mathrm{alg}} = 2960{,}39\ \mathrm{К}$, а логарифмическая проекция Саха — $T_{\mathrm{rec}} = 2977{,}35\ \mathrm{К}$ (наблюдение Planck~2018: $T_* = 2972{,}98\ \mathrm{К}$, отклонение $0{,}42\%$ и~$0{,}15\%$).
\item \textbf{Коэффициент космологического остывания $1 + z_{\mathrm{CMB}}$ и современный квант $\Theta_{\mathrm{CMB}}$:}
  при расширении пузыря от такта~$N_{\mathrm{CMB}}$ до такта~$N_{\mathrm{today}}$ в размерности~$d$ тактовый фактор растяжения совпадает с отношением ступеней башни~$\vevHierarchySteps$ к~$\alpha_{\mathrm{fs}}$:
  \begin{equation}
  \label{eq:cosmo-redshift-dilation}
    1 + z_{\mathrm{CMB}}
    = \left(\frac{N_{\mathrm{today}}}{N_{\mathrm{CMB}}}\right)^{(d-1)/d}
    \approx 1095{,}95
    \;\longleftrightarrow\;
    \frac{\vevHierarchySteps}{\alpha_{\mathrm{fs}}}
    \approx 1096{,}29
  \end{equation}
  (совпадение с точностью~$0{,}031\%$; наблюдение Planck~2018: $1 + z_* = 1090{,}80 \pm 0{,}21$).
  Отсюда современный тепловой квант реликтового излучения в единицах~$E_{\mathrm{Ry}}$, $m_e c^2$ и сверхтонкого кванта водорода~$E_{\mathrm{HI}}$ равен:
  \begin{equation}
  \label{eq:cosmo-theta-cmb-natural}
  \begin{aligned}
    \frac{\Theta_{\mathrm{CMB}}}{E_{\mathrm{Ry}}}
    &= \frac{d\,\alpha_{\mathrm{fs}}}{d_{\mathrm{in}}(\square)\,d_{\mathrm{out}}(E)\,\vevHierarchySteps^{d-1}},\\
    \frac{\Theta_{\mathrm{CMB}}}{m_e c^2}
    &= \frac{d\,\alpha_{\mathrm{fs}}^{d}}{(d-1)\,d_{\mathrm{in}}(\square)\,d_{\mathrm{out}}(E)\,\vevHierarchySteps^{d-1}},
    \qquad
    \frac{\Theta_{\mathrm{CMB}}}{E_{\mathrm{HI}}}
    \approx d_{\mathrm{out}}(E)\,\vevHierarchySteps.
  \end{aligned}
  \end{equation}
\item \textbf{Запись в собственной единице такта носителя $T_{P,\Mlayer} := E_0 = s_0/\htick$ через чисто геометрический инвариант $r_{\mathrm{disc}} = \Delta\varphi_{\mathrm{disc}}/\phaseRingCount$ (без $\pi$ и без $\ln 2$):}
  все целочисленные параметры регистра ячейки задаются непосредственно комбинаторикой граней и каналов однотактового кубооктаэдра ($|\square|=2d=6$, $|\triangle|=2^d=8$, $|\mathrm{N12}|=12$, $N_\varphi=1+|\mathrm{N12}|=13$, $d_{\mathrm{in}}(\square)=d+1=4$, $d_{\mathrm{out}}(E)=5$):
  \[
  \begin{aligned}
    \mathrm{frac\_bits} &= |\square| = 6, &
    \vevHierarchySteps &= |\triangle| = 8,\\
    \mathrm{mod\_bits} &= 1 + |\triangle| = d^2 = 9, &
    \phaseRingCount &= |\triangle|^d = 2^{1+|\triangle|} = 512,
  \end{aligned}
  \]
  а порог Гейзенберга $\Delta\varphi_{\mathrm{disc}} = 41$ определяется четырьмя тождественно совпадающими геометрическими соотношениями без участия~$\pi$:
  \begin{equation}
  \label{eq:cosmo-rdisc-geometric}
  \begin{aligned}
    \Delta\varphi_{\mathrm{disc}}
    &= 1 + d_{\mathrm{out}}(E)\,|\triangle|
    = d_{\mathrm{in}}(\square)^2 + d_{\mathrm{out}}(E)^2
    = N_{\mathrm{shell}}(2) - 1
    = \frac{\phaseRingCount + N_\varphi + |\triangle|}{N_\varphi}
    = 41,\\
    r_{\mathrm{disc}}
    &= \frac{\Delta\varphi_{\mathrm{disc}}}{\phaseRingCount}
    = \frac{1 + d_{\mathrm{out}}(E)\,|\triangle|}{|\triangle|^d}
    = \frac{41}{512} \approx 0{,}0800781
  \end{aligned}
  \end{equation}
  (тогда как континуальная величина $1/(4\pi) \approx 0{,}0795775$, отличающаяся на фиксированный сдвиг $+0{,}63\%$, представляет собой лишь длинноволновое приближение~$r_{\mathrm{disc}}$ при сферическом усреднении на~$\Tlayer$).
  Подстановка $r_{\mathrm{disc}}$ ($v/E_0 = \alpha_{\mathrm{fs}}^{\vevHierarchySteps}/\sqrt{r_{\mathrm{disc}}}$, $\lambda_H = 1/\vevHierarchySteps + \vevHierarchySteps r_{\mathrm{disc}}\alpha_{\mathrm{fs}}$, $m_H/E_0 = \sqrt{(d-1)\lambda_H}\,v/E_0$ и $E_{\mathrm{Ry}}/E_0 = \frac{\alpha_{\mathrm{fs}}^{|\mathrm{N12}|}}{(d-1)N_\varphi}\sqrt{(d-1)\lambda_H/r_{\mathrm{disc}}}$) выражает оба космологических тепловых кванта через точные структурные показатели $|\mathrm{N12}|$ и $N_\varphi = |\mathrm{N12}| + 1$:
  \begin{equation}
  \label{eq:cosmo-theta-cmb-in-e0}
  \begin{aligned}
    \frac{\Theta_{\mathrm{rec}}}{E_0}
    &= \frac{d\,\alpha_{\mathrm{fs}}^{|\mathrm{N12}|}}{(d-1)\,d_{\mathrm{in}}(\square)\,d_{\mathrm{out}}(E)\,\vevHierarchySteps\,N_\varphi}\sqrt{\frac{(d-1)\lambda_H}{r_{\mathrm{disc}}}}
    \approx 2{,}959\times 10^{-29},\\
    \frac{\Theta_{\mathrm{CMB}}}{E_0}
    &= \frac{d\,\alpha_{\mathrm{fs}}^{N_\varphi}}{(d-1)\,d_{\mathrm{in}}(\square)\,d_{\mathrm{out}}(E)\,\vevHierarchySteps^{d}\,N_\varphi}\sqrt{\frac{(d-1)\lambda_H}{r_{\mathrm{disc}}}}
    \approx 2{,}699\times 10^{-32}.
  \end{aligned}
  \end{equation}
  На приборном слое~$\Tlayer$ эти соотношения дают $T_{\mathrm{CMB}}^{\mathrm{alg}} = 2{,}7004\ \mathrm{К}$ и $T_{\mathrm{CMB}}^{\mathrm{Saha}} = 2{,}7158\ \mathrm{К}$ (эталон COBE/FIRAS: $2{,}72548 \pm 0{,}00057\ \mathrm{К}$, отклонение $0{,}35\%\text{--}0{,}92\%$).
\end{enumerate}
\end{theorem}

\begin{lemma}[{О пятёрке расщепляющихся простых чисел $\{5, 13, 41, 97, 137\}\subset\mathbb{Z}$ как нормах гауссовых простых в $\mathbb{Z}[i]$}]
\label{lem:cosmo-five-splitting-primes-zi}
Поскольку спинорный регистр носителя~$\Mlayer$ построен над кольцом гауссовых целых чисел $\mathbb{Z}[i]$ ($\mathbb{Z}_{\phaseRingCount}[i]^2$), фундаментальные простые инварианты кубооктаэдра разделяются по закону взаимности Ферма--Эйлера: инертные простые вида $q \equiv 3 \pmod{d_{\mathrm{in}}(\square)}$ ($d = d_{\mathrm{in}}(\triangle) = 3$, $d_{\mathrm{out}}(V) = d_{\mathrm{soft}} = 7$, $|\mathrm{N12}| - 1 = 11$) сохраняют простоту в~$\mathbb{Z}[i]$, тогда как пятёрка ключевых простых знаменателей и порогов лестницы
\[
  \bigl\{d_{\mathrm{out}}(E),\; N_\varphi,\; \Delta\varphi_{\mathrm{disc}},\; \hat M,\; \alpha_{\mathrm{geom}}^{-1}\bigr\}
  = \bigl\{5,\; 13,\; 41,\; 97,\; 137\bigr\} \subset \mathbb{Z}
\]
состоит из рациональных простых вида $p \equiv 1 \pmod{d_{\mathrm{in}}(\square)}$, каждое из которых расщепляется в~$\mathbb{Z}[i]$ в произведение пары сопряжённых гауссовых простых $p = \pi_p\bar\pi_p = N_{\mathbb{Z}[i]/\mathbb{Z}}(\pi_p)$, порождённых внутренними степенями четырёх каналов кубооктаэдра $d_{\mathrm{in}}(V,E,\triangle,\square) = (1,2,3,4)$:
\begin{equation}
\label{eq:cosmo-five-splitting-primes-zi}
\begin{aligned}
  d_{\mathrm{out}}(E) &= 5 = (1+2i)(1-2i) = d_{\mathrm{in}}(V)^2 + d_{\mathrm{in}}(E)^2,\\
  N_\varphi &= 13 = (2+3i)(2-3i) = d_{\mathrm{in}}(E)^2 + d_{\mathrm{in}}(\triangle)^2,\\
  \Delta\varphi_{\mathrm{disc}} &= 41 = (4+5i)(4-5i) = d_{\mathrm{in}}(\square)^2 + d_{\mathrm{out}}(E)^2,\\
  \hat M &= 97 = (4+9i)(4-9i) = d_{\mathrm{in}}(\square)^2 + \mathrm{mod\_bits}^2,\\
  \alpha_{\mathrm{geom}}^{-1} &= 137 = (4+11i)(4-11i) = d_{\mathrm{in}}(\square)^2 + \bigl(|\mathrm{N12}| - 1\bigr)^2.
\end{aligned}
\end{equation}
При этом два способа умножения первых двух гауссовых простых $\pi_5 = 1 + 2i$ и $\pi_{13} = 2 \pm 3i$ в~$\mathbb{Z}[i]$ ($\pi_5\pi_{13} = -4 + 7i$ и $\pi_5\bar\pi_{13} = 8 + i$) тождественно связывают знаменатель рекомбинации $d_{\mathrm{out}}(E)\,N_\varphi = 65$ с внешним вершинным каналом $d_{\mathrm{out}}(V) = 7$ и высотой башни $\vevHierarchySteps = |\triangle| = 8$:
\begin{equation}
\label{eq:cosmo-brahmagupta-65}
  d_{\mathrm{out}}(E)\,N_\varphi
  = d_{\mathrm{in}}(\square)^2 + d_{\mathrm{out}}(V)^2
  = \vevHierarchySteps^2 + d_{\mathrm{in}}(V)^2
  = 65.
\end{equation}
\end{lemma}

\section{Двухкаскадная эволюция сида \texorpdfstring{$\mathrm{IV}=263$}{IV=263} через полупериодное горлышко \texorpdfstring{$\tau_{\mathrm{neck}} = \phaseRingCount/2 = 256$}{t=256} и природа ранних ядер JWST}
\label{sec:cosmo-two-epoch-bottleneck}

Длинный побитово-точный прогон канонического 9-битного сида $\mathrm{IV} = 263$ на трёхмерной решётке~\latGCKshort\ (при фиксированном фазовом классе центра $\mathrm{center\_class} = 11$ траектория вплоть до $\tau = 512$ идентична на кубах $36^3$, $48^3$, $54^3$ и $60^3$ за счёт гашения подпороговой периферии затвором Гейзенберга $\Delta\varphi_{\mathrm{disc}} = 41$) выявляет двухкаскадную структуру рождения материи и излучения (CE-C-04):
\begin{enumerate}[leftmargin=*]
\item \textbf{Первичный фотонный шар на внутреннем горлышке $\tau_1 = \vevHierarchySteps = 8$:}
  на такте $\tau = 8$ число барионных вихрей обращается в ноль ($b_{\mathrm{hits}} = 0$), тогда как надпороговое фазовое поле содержит $1307$ фотонных квантов Гейзенберга ($n_E = \lfloor|\Phi|/41\rfloor \ge 1$) при максимальном контрасте плотности $\rho_{\max}/\langle\rho\rangle = 256{,}8$. Из этого первичного фотонного шара на тактах $\tau = 16\ldots 64$ конденсируется первая волна из $b_{\mathrm{hits}} = 6$ единичных барионных вихрей ($w_{\max} = 1$).
\item \textbf{Полупериодное горлышко отцепления $\tau_{\mathrm{neck}} = \phaseRingCount/2 = 256$ (срез реликтового излучения и одиночного зерна):}
  к полупериоду фазового кольца $\tau = 256$ пять из шести вихрей первой волны сбрасывают фазу в излучение. Каузальный пузырь ($1047$ активных ячеек радиуса $R_{\mathrm{causal}} \approx 9{,}06$) заполняется однородной фотонной баней из $2733$ квантов Гейзенберга (с отделившимся внешним фотонным гало из $65$ квантов в безбарионной оболочке $r \in [8, 14)$ и стабильностью средней плотности $\Delta\langle\rho\rangle/\langle\rho\rangle < 0{,}2\%$). Посреди этой фотонной бани выживает ровно один локализованный барионный реликт ($b_{\mathrm{hits}} = 1$, $w = -1$) со смещением $\Delta\mathbf{x}_{\mathrm{relic}} = (-2, -3, +2)$ от центра пузыря ($r = \sqrt{17} \approx 4{,}12$).
\item \textbf{Рождение трёхквантового ядра ($w_{\max} = 3 = d$) на соседнем ребре~\latGCKshort\ во Второй Эпохе и следствие для JWST:}
  после прохождения горлышка $\tau_{\mathrm{neck}} = 256$ выживший реликтовый узел служит готовым центром конденсации для окружающей фотонной бани. На такте $\tau = 336$ буквально в одном шаге ближайшего соседа решётки~$N_{12}$ от реликтового зерна — в узле со смещением $(-1, -3, +3) = \Delta\mathbf{x}_{\mathrm{relic}} + (+1, 0, +1)$ — формируется тяжёлое трёхквантовое ядро с индексом намотки $|w| = 3 = d$ (и парное двухквантовое ядро $|w| = 2$ на втором соседнем ребре $\Delta\mathbf{x}_{\mathrm{relic}} + (-1, +1, 0)$), после чего на такте $\tau = 352$ вспыхивает кластер из $b_{\mathrm{hits}} = 8 = \vevHierarchySteps$ барионных узлов.
  В макроскопической космологии это объясняет одновременное выполнение трёх наблюдаемых фактов:
  \begin{enumerate}[label=(\alph*),noitemsep]
  \item отсутствие крупномасштабных анизотропных колец в карте Planck CMB (поскольку к моменту $\tau_{\mathrm{neck}}$ фотонная баня полностью термализована с вариацией $< 0{,}2\%$);
  \item раннее появление зрелых массивных ядер и галактик, регистрируемых телескопом JWST на $z \sim 10\text{--}15$ (поскольку конденсация Второй Эпохи идёт непосредственно на соседних рёбрах~$N_{12}$ вокруг выживших с Первой Эпохи точечных топологических зёрен);
  \item различие между видимым возрастом разлёта текущей барионной волны после горлышка ($\tau_{\mathrm{wave}} = \tau - \tau_{\mathrm{neck}} \leftrightarrow 13{,}8\ \mathrm{млрд\ лет}$) и полным каузальным возрастом пузыря от момента нуклеации сида $\mathrm{IV}=263$ ($\tau_{\mathrm{total}} = \tau_{\mathrm{neck}} + \tau_{\mathrm{wave}} \sim 26\text{--}27\ \mathrm{млрд\ лет}$).
  \end{enumerate}
\item \textbf{Первичный нуклеосинтез (BBN: массовая доля гелия-4~$Y_p$ и отношение $n_{\mathrm{D}}/n_{\mathrm{H}}$) и перепись намоток $|w|\in\{1,2,3\}$ во Второй Эпохе:}
  на пороге замораживания слабых переходов отношение числа нейтронов к протонам равно $n/p = 1/d_{\mathrm{soft}} = 1/d_{\mathrm{out}}(V)$ (один канал из $d_{\mathrm{soft}} = \vevHierarchySteps - 1 = 7$ вершинных каналов), откуда древесная массовая доля первичного гелия-4 составляет $Y_p^{(0)} = \frac{(d-1)(n/p)}{1 + n/p} = \frac{d-1}{\vevHierarchySteps} = \frac{1}{d_{\mathrm{in}}(\square)} = 0{,}25$, а с учётом распада свободных нейтронов на шве кластера $d_{\mathrm{in}}(\square)N_\varphi$ и квадратичного ветвления по $d_{\mathrm{out}}(E)=5$ внешним рёберным каналам:
  \begin{equation}
  \label{eq:cosmo-bbn-yp-dh}
    Y_p = \frac{1}{d_{\mathrm{in}}(\square)}\!\left(1 - \frac{1}{d_{\mathrm{in}}(\square)\,N_\varphi}\right) \approx 0{,}24519,
    \qquad
    \frac{n_{\mathrm{D}}}{n_{\mathrm{H}}} = \frac{d_{\mathrm{out}}(E)^{d-1}}{d_{\mathrm{in}}(\square)\,N_\varphi}\,\alpha_{\mathrm{fs}}^{d-1} \approx 2{,}560\times 10^{-5}
  \end{equation}
  (наблюдения PDG / Planck~2018: $Y_p = 0{,}245 \pm 0{,}003$, откл.~$0{,}08\%$; $n_{\mathrm{D}}/n_{\mathrm{H}} = (2{,}547 \pm 0{,}025)\times 10^{-5}$, откл.~$0{,}52\%$).
  Прямой подсчёт вихревых узлов по индексам намотки $|w|\in\{1,2,3\}$ во Второй Эпохе прогона $\mathrm{IV}=263$ ($\tau = 272\ldots 512$ с шагом $(d-1)\vevHierarchySteps = 16$) фиксирует $61$ одноквантовый узел $|w|=1$, $11$ двухквантовых узлов $|w|=2$ и $1$ трёхквантовый узел $|w|=d$. На этой конечной выборке доля суммарной массы намотки $\sum |w|\,N_{|w|}$, приходящаяся на двухквантовый гелиевый канал $|w|=2$, равна $22/86 \approx 0{,}2558$, что совпадает с древесным порогом $1/d_{\mathrm{in}}(\square) = 0{,}2500$ с точностью $|22/86 - 1/4| = 1/172 \approx 0{,}0058$, то есть глубоко внутри одного дискретного кванта вихревой массы $\Delta w / \sum |w|N_{|w|} = 2/86 \approx 0{,}0233$ (и пуассоновского сдвига $\pm 1$ узел $\approx \pm 0{,}017$).
\end{enumerate}

\input{volumes/include-computational-laboratory.tex}

\input{volumes/include-measurement-processing.tex}

\chapter{Критические эксперименты и фальсификация космологического тома}
\label{ch:cosmo-critical-experiments}

Журнал критических вычислительных опытов космологического тома (серия CE-C: нуклеация барионной асимметрии из $\mathrm{IV}=263$, кубооктаэдрический ореол $16/3$ и строгая upstream-сверка $\eta_B$ и $\Lambda_{\mathrm{bubble}}$ по протоколу GUM главы~\ref{ch:measurement-processing}):

% Журнал CE-C — том cosmology.

\begin{labseries}{Серия CE-C: нуклеация из $\mathrm{IV}=263$, ореол $16/3$ и космологические инварианты}{sec:lab-series-ce-c}

\begin{labrecord}{CE-C-01 · Нуклеация барионной асимметрии из 9-битного сида $\mathrm{IV}=263$ против $\mathrm{IV}=7$}{lab:ce-c-01}
\label{exp:cec01}

\subsubsection*{Постановка}
\textbf{Физический вопрос.}
На микроуровне~$\Mlayer$ локальный закон~$g$ строго обратим ($A_{13}$) и симметричен относительно обращения знака фазы, однако наш каузальный пузырь на макроуровне~$\Tlayer$ содержит избыток барионов ($n=+1$) над антибарионами ($n=-1$).
Способен ли детерминированный шаг~$g$ на трёхмерной решётке~\latGCKshort\ без посаженных начальных вихрей ($b(\tau=0)\equiv 0$) породить устойчивые барионные солитоны ($b=1$) из 9-битного начального вектора зародыша $\mathrm{IV}_{\mathrm{bubble}} = 2^{\vevHierarchySteps} + d_{\mathrm{soft}} = 263 = \texttt{0b1\_00000111}$ (\cref{thm:cosmo-9bit-iv-baryon}), и меняет ли знак рождающегося мира инверсия старшего бита хиральности ($\mathrm{IV} = 7 = \texttt{0b0\_00000111}$)?

\textbf{Что проверяем.}
\begin{enumerate}[label=\textbf{(\alph*)},leftmargin=*,itemsep=0.25em]
\item Прямой прогон 3D-решётки~\latGCKshort\ при $\mathrm{IV} = 263$ ($b_{\mathrm{chir}}=1$, $\mathrm{mask}_8 = 7$): нуклеация вихревых ядер ($b=1$, \texttt{born=1}) и положительный суммарный топологический заряд $\sum n > 0$.
\item Зеркальный контроль $\mathrm{IV} = 7$ ($b_{\mathrm{chir}}=0$, $\mathrm{mask}_8 = 7$): инверсия знака заряда ($\sum n < 0$, антиматерия).
\item Нулевые контроли: однородный сдвиг ($\mathrm{mask}_8 = 0$) и плоская одномерная стенка (\texttt{wall}, $\mathrm{mask}_8 = 15$), где контурный интеграл фазы равен нулю ($\oint\mathrm{d}\varphi = 0 \implies \sum b \equiv 0$).
\end{enumerate}

\textbf{Рабочая гипотеза.}
Только трёхгранный узел Киббла--Зурека на стыке трёх ортогональных октантов барионного треугольника~$K_3$ ($\mathrm{mask}_8 = 7 = d_{\mathrm{soft}}$) создаёт ненулевую голономию и переводит фазовый импульс через порог дискретности $\Delta\varphi_{\mathrm{disc}} = 41$, а знак $\sum n$ однозначно определяется старшим битом~$b_{\mathrm{chir}}$.

\textbf{Наблюдаемые.}
Число родившихся ядер $N_{\mathrm{born}} = \sum_x b(x)$, суммарный топологический заряд $Q_{\mathrm{tot}} = \sum_x n(x)$ и булев индикатор нуклеации \texttt{born} на такте $\tau = 32$ (табл.~\ref{tab:exp-ce-c-01}).

\textbf{Критерий согласия с теорией.}
Детерминированное выполнение $\texttt{born}=1$, $Q_{\mathrm{tot}} > 0$ при $\mathrm{IV}=263$; $Q_{\mathrm{tot}} < 0$ при $\mathrm{IV}=7$; строгое тождество $N_{\mathrm{born}} \equiv 0$ при $\mathrm{mask}_8 \in \{0, 15\}$ (тип~B, целочисленная арифметика).

\subsubsection*{Теоретическая часть}
По теореме~\ref{thm:cosmo-9bit-iv-baryon} три координатные плоскости вокруг центра нуклеации делят пространство на $2^d = 8 = \vevHierarchySteps$ октантов, что в сумме с битом хиральности спинора $b_{\mathrm{chir}}\in\{0,1\}$ исчерпывает весь 9-битный регистр ячейки Бекештейна ($\dim_2(\mathrm{IV}) = 1 + 8 = 9 = \lfloor\caBitBudget\rfloor$).
Плоская стенка меняет фазу лишь вдоль одной оси и даёт $\oint_{\partial f}\Delta\varphi = 0$ на каждой треугольной плакетке~$K_3$, тогда как трёхгранный угол $\mathrm{mask}_8 = 7$ концентрирует скачок фазы на общей вершине трёх плоскостей $\{XY, YZ, XZ\}$.

\textbf{Проверка размерностей.}
Все входы и выходы опыта строго безразмерны:
\[
  [\mathrm{IV}] = 1\quad(\mathbb{Z}_{512}),\qquad
  [\tau/\htick] = 1,\qquad
  [b(x)] = 1\quad(\{0,1\}),\qquad
  \left[\sum_x n(x)\right] = 1\quad(\mathbb{Z}),
\]
а перевод к кулонам в СИ осуществляется умножением целого~$n$ на элементарный заряд~$e_0$ ($[Q_{\mathrm{SI}}] = [n\,e_0] = \mathrm{I\cdot T}$).

\subsubsection*{Ход эксперимента}
Скрипт воспроизведения: \texttt{python scripts/run\_phase\_iv\_to\_i.py --fcc}
(функция генерации сида в \texttt{mt\_ca/seeds.py}).

\begin{enumerate}[label=\textbf{Шаг \arabic*:},leftmargin=*,itemsep=0.35em]
\item Инициализировать 3D-решётку~\latGCKshort\ ($16^3$) в режиме кипящего пола~$A_5$ с $b(\tau=0)\equiv 0$.
\item Подать 9-битный вектор $\mathrm{IV}\in\{263, 7, 0, 15\}$ и выполнить $32$ обратимых шага~$g$.
\item Снять $N_{\mathrm{born}} = \sum_x b(x)$ и $Q_{\mathrm{tot}} = \sum_x n(x)$ в табл.~\ref{tab:exp-ce-c-01}.
\end{enumerate}

\begin{table}[htbp]
\centering
\footnotesize
\setlength{\tabcolsep}{4pt}
\caption{Нуклеация барионной асимметрии из 9-битного вектора $\mathrm{IV}\in\mathbb{Z}_{512}$ на 3D~\latGCKshort\ (CE-C-01)}
\label{tab:exp-ce-c-01}
\begin{tabular}{@{}llllr@{}}
\toprule
Режим сида & $\mathrm{IV}\in\mathbb{Z}_{512}$ & $b_{\mathrm{chir}}$ & Маска октантов $\mathrm{mask}_8$ & Итог к $\tau=32$ ($\texttt{born}$, $\sum n$) \\
\midrule
Наш пузырь ($K_3$, материя) & $263$ (\texttt{0b1\_00000111}) & $1$ & $7 = d_{\mathrm{soft}}$ ($3$ октанта $K_3$) & $\texttt{born}=1,\ \sum n = +2 > 0$ \\
Зеркальный (антиматерия) & $7$ (\texttt{0b0\_00000111}) & $0$ & $7 = d_{\mathrm{soft}}$ ($3$ октанта $K_3$) & $\texttt{born}=1,\ \sum n = -2 < 0$ \\
Однородный лёд (контроль) & $0$ (\texttt{0b0\_00000000}) & $0$ & $0$ (без градиента) & $\texttt{born}=0,\ \sum b \equiv 0$ \\
Плоская стенка (\texttt{wall}) & $15$ (\texttt{0b0\_00001111}) & $0$ & полупространство $x\ge c_x$ & $\texttt{born}=0,\ \sum b \equiv 0$ \\
\bottomrule
\end{tabular}
\end{table}

\subsubsection*{Анализ, расчёт погрешности и выводы}
Поскольку ядро работает в точной модулярной арифметике $\mathbb{Z}_{512}^4$, статистическая погрешность типа~A равна нулю ($u_A = 0$), а погрешность типа~B отсутствует на целочисленных инвариантах ($b\in\{0,1\}$, $n\in\mathbb{Z}$).
Таблица~\ref{tab:exp-ce-c-01} доказывает, что трёхгранный узел $\mathrm{IV}=263$ детерминированно рождает материю ($\sum n = +2$), а замена старшего бита $\mathrm{IV}=7$ зеркально рождает антиматерию ($\sum n = -2$).

\textbf{Вывод.}
Барионная асимметрия нашего пузыря однозначно кодируется 9-битным сидом $\mathrm{IV}=263\in\mathbb{Z}_{512}$ без нарушения обратимости закона~$g$.
\textbf{Статус записи:} опыт закрыт по табл.~\ref{tab:exp-ce-c-01}.

\end{labrecord}

\begin{labrecord}{CE-C-02 · Отношение нейтрального ореола кубооктаэдра к ядру Вороного $\Omega_{\mathrm{DM}}/\Omega_b = 16/3$}{lab:ce-c-02}
\label{exp:cec02}

\subsubsection*{Постановка}
\textbf{Физический вопрос.}
В астрофизических наблюдениях (Planck~2018) отношение плотности гравитационного ореола («тёмной материи») к плотности барионного вещества составляет $\Omega_{\mathrm{DM}}/\Omega_b = 5{,}364 \pm 0{,}060$.
Воспроизводится ли это число без введения новых частиц как геометрическое отношение объёма однотактового кубооктаэдра упругой деформации $V_{\mathrm{cuboct}}$ к объёму элементарной ячейки Вороного локализованного барионного ядра $dV_{\mathrm{Voronoi}}$ на решётке~\latGCKshort\ (\cref{prop:cosmo-dm-baryon-ratio})?

\textbf{Что проверяем.}
\begin{enumerate}[label=\textbf{(\alph*)},leftmargin=*,itemsep=0.25em]
\item Тождество полноты носителя $D(x) - (D_{\mathrm{baryon}}(\mu) - \delta K(x)) \equiv 0$ (нулевой «тёмный остаток» $D_{\mathrm{dark}}\equiv 0$), нетривиальный вакуумный импеданс $D_{\mathrm{vac}}(x) = 9856 - \sum_{O_2}|A_{xy}|$ на единичном полу вакуума $\mu_{\mathrm{vac}}=1$ и точное целочисленное отношение ёмкости ореола кубооктаэдра к $3\mathrm{D}$-ядру Вороного $W_{\mathrm{halo}}/W_{\mathrm{core}} = 2048/384 = V_{\mathrm{cuboct}}/dV_{\mathrm{Voronoi}} = 2\vevHierarchySteps/d = 16/3$ (\texttt{no\_\allowbreak dark\_\allowbreak matter\_\allowbreak and\_\allowbreak vacuum\_\allowbreak holonomy\_\allowbreak halo} в \texttt{UnifiedOperators.lean}).
\item Согласие модельных отношений $\Omega_{\mathrm{DM}}/\Omega_b = 16/3$, $\Omega_m/\Omega_b = 19/3$ и барионной доли $\Omega_b/\Omega_m = 3/19$ с данными Planck~2018 по критерию $|E_n| \le 1$ (\cref{ch:measurement-processing}).
\end{enumerate}

\textbf{Рабочая гипотеза.}
На гранецентрированной кубической решётке отношение объёма первой каузальной оболочки ($\mathrm{strain}$-ореола) к объёму узла Вороного равно $16/3 \approx 5{,}3333$, что лежит внутри $1\sigma$-интервала Planck~2018 ($|E_n| < 0{,}5$).

\textbf{Наблюдаемые.}
$\Omega_{\mathrm{DM}}/\Omega_b$, $\Omega_m/\Omega_b$, $\Omega_b/\Omega_m$ и их нормированные отклонения~$|E_n|$ (табл.~\ref{tab:exp-ce-c-02}).

\textbf{Критерий согласия с теорией.}
Алгебраическое равенство $V_{\mathrm{cuboct}}/dV_{\mathrm{Voronoi}} = 16/3$ (остаток $<10^{-15}$) и показатель $|E_n| \le 1$ по всем трём отношениям (тип~B).

\subsubsection*{Теоретическая часть}
По предложению~\ref{prop:cosmo-dm-baryon-ratio} барионное ядро ($b=1$) занимает ромбический додекаэдр Вороного объёма $dV_{\mathrm{Voronoi}} = \caKappaGeom a^3 = a^3/\sqrt{2}$, тогда как его однотактовая каузальная оболочка натяжения Редже охватывает кубооктаэдр объёма $V_{\mathrm{cuboct}} = \frac{2\vevHierarchySteps}{d}\caKappaGeom a^3 = \frac{8\sqrt{2}}{3}a^3$.

\textbf{Проверка размерностей.}
Объёмы числителя и знаменателя имеют размерность куба длины $\mathrm{L}^3$, поэтому их отношение строго безразмерно и инвариантно относительно масштаба~$a$:
\[
  [V_{\mathrm{cuboct}}] = [dV_{\mathrm{Voronoi}}] = [a^3] = \mathrm{L^3}
  \quad\Longrightarrow\quad
  \left[\frac{\Omega_{\mathrm{DM}}}{\Omega_b}\right] = \left[\frac{V_{\mathrm{cuboct}}}{dV_{\mathrm{Voronoi}}}\right] = \frac{\mathrm{L^3}}{\mathrm{L^3}} = 1.
\]

\subsubsection*{Ход эксперимента}
Воспроизведение: \texttt{python verify\_principles.py --suite sm} (модуль \texttt{mt\_ca/si\_sm\_rows.py}).

\begin{enumerate}[label=\textbf{Шаг \arabic*:},leftmargin=*,itemsep=0.35em]
\item Вычислить $dV_{\mathrm{Voronoi}}/a^3 = \caKappaGeom$ и $V_{\mathrm{cuboct}}/a^3 = (2\vevHierarchySteps/d)\caKappaGeom$ из констант носителя.
\item Сопоставить отношения $16/3$, $19/3$ и $3/19$ с таблицей Planck~2018 ($\Omega_c h^2 = 0{,}1200\pm 0{,}0012$, $\Omega_b h^2 = 0{,}02237\pm 0{,}00015$).
\item Рассчитать $|E_n| = |x_{\mathrm{model}} - x_{\mathrm{ref}}| / (2 u_{\mathrm{ref}})$ при коэффициенте охвата $k=2$.
\end{enumerate}

\begin{table}[htbp]
\centering
\footnotesize
\setlength{\tabcolsep}{4pt}
\caption{Отношение нейтрального ореола кубооктаэдра к барионному ядру Вороного (CE-C-02)}
\label{tab:exp-ce-c-02}
\begin{tabular}{@{}lrrrr@{}}
\toprule
Наблюдаемое & Формула модели & Модель & Planck 2018 ($1\sigma$) & $|E_n|$ ($k=2$) \\
\midrule
$\Omega_{\mathrm{DM}}/\Omega_b$ & $V_{\mathrm{cuboct}}/dV_{\mathrm{Voronoi}} = 2\vevHierarchySteps/d = 16/3$ & $5{,}3333$ & $5{,}364 \pm 0{,}060$ & $0{,}26$ \\
$\Omega_m/\Omega_b$ & $1 + V_{\mathrm{cuboct}}/dV_{\mathrm{Voronoi}} = 19/3$ & $6{,}3333$ & $6{,}364 \pm 0{,}060$ & $0{,}26$ \\
$\Omega_b/\Omega_m$ & $d/(2\vevHierarchySteps + d) = 3/19$ & $0{,}15789$ & $0{,}1571 \pm 0{,}0015$ & $0{,}26$ \\
\bottomrule
\end{tabular}
\end{table}

\subsubsection*{Анализ, расчёт погрешности и выводы}
Теоретическое отношение $16/3$ является точной геометрической константой ($u_{\mathrm{model}} = 0$); вся неопределённость сравнения определяется погрешностью астрофизического эталона Planck~2018 (тип~B, $u_{\mathrm{ref}} = 0{,}060$).
Значение $|E_n| = 0{,}26 \ll 1$ означает полное метрологическое согласие на уровне $0{,}5\sigma$.

\textbf{Вывод.}
Эффект «тёмной материи» количественно исчерпывается гравитационным вкладом однотактового кубооктаэдра деформации вокруг барионной ячейки Вороного ($16/3$).
\textbf{Статус записи:} опыт закрыт по табл.~\ref{tab:exp-ce-c-02}.

\end{labrecord}

\begin{labrecord}{CE-C-03 · Строгая upstream-сверка $\eta_B$, $N_{\mathrm{today}}$ и горизонтной кривизны $\Lambda_{\mathrm{bubble}}$}{lab:ce-c-03}
\label{exp:cec03}

\subsubsection*{Постановка}
\textbf{Физический вопрос.}
Замыкаются ли в единую некруговую цепочку от постоянной тонкой структуры~$\alpha_{\mathrm{fs}}$ (\bcref{thm:alpha}) и такта $\htick = \caKappaGeom t_P$ (\cref{thm:cosmo-time-bridge}) три главных космологических инварианта нашего пузыря: барион-фотонное отношение $\eta_B = n_B/n_\gamma$ (\cref{thm:cosmo-9bit-iv-baryon}), число тактов возраста $N_{\mathrm{today}}$ и горизонтная кривизна $\Lambda_{\mathrm{bubble}}\ell_P^2$ со сферическим коэффициентом проектора~$\mathcal{B}$ (\cref{sec:cosmo-lambda-solution})?

\textbf{Что проверяем.}
\begin{enumerate}[label=\textbf{(\alph*)},leftmargin=*,itemsep=0.25em]
\item Строгий upstream-вывод $\eta_B = \bar\eta\,(d_{\mathrm{out}}(E)/\vevHierarchySteps)\,\alpha_{\mathrm{fs}}^4 = \frac{45}{208}\alpha_{\mathrm{fs}}^4$ без подстановки астрофизических подгонок.
\item Точный перевод возраста пузыря $t_{\mathrm{age}} \approx 13{,}787\ \mathrm{млрд\ лет}$ в целое число тактов $N_{\mathrm{today}} = \lfloor t_{\mathrm{age}}/\htick\rfloor$.
\item Устранение зазора между кубическим потолком $|\square|/N_{\mathrm{today}}^2 \approx 4{,}61\times 10^{-122}$ и наблюдением Planck~2018 ($2{,}89\times 10^{-122}$) за счёт геометрического коэффициента сферической проекции $\eta_{\mathrm{surf}} = A_{\mathrm{sphere}}/A_{\mathrm{cuboct}} = \pi/(3+\sqrt{3}) \approx 0{,}6639$.
\end{enumerate}

\textbf{Рабочая гипотеза.}
Учёт поверхностного фактора проектора~$\mathcal{B}$ ($\eta_{\mathrm{surf}} \approx 0{,}6639$) переводит $\Lambda_{\mathrm{bubble}}\ell_P^2$ из $4{,}61\times 10^{-122}$ в $3{,}06\times 10^{-122}$, сводя нормированное отклонение~$|E_n|$ как по~$\eta_B$, так и по~$\Lambda_{\mathrm{bubble}}$ внутрь единичного порога $|E_n| < 1$.

\textbf{Наблюдаемые.}
$\eta_B$, $N_{\mathrm{today}}$, $\Lambda_{\mathrm{bubble}}\ell_P^2$ и их показатели $|E_n|$ относительно Planck~2018 (табл.~\ref{tab:exp-ce-c-03}).

\textbf{Критерий согласия с теорией.}
Показатель $|E_n| \le 1$ при $k=2$ для всех трёх наблюдаемых без свободных параметров (тип~B, протокол главы~\ref{ch:measurement-processing}).

\subsubsection*{Теоретическая часть}
По теореме~\ref{thm:cosmo-9bit-iv-baryon} барионная асимметрия возникает в четвёртом порядке по~$\alpha_{\mathrm{fs}}$ с геометрическим множителем $\bar\eta\cdot d_{\mathrm{out}}(E)/\vevHierarchySteps = (9/26)\cdot(5/8) = 45/208$.
Согласно~\ref{sec:cosmo-lambda-solution} (\eqref{eq:cosmo-lambda-bubble}), макроскопический горизонт радиуса $R_H = \caKappaGeom N_{\mathrm{today}}\ell_P$ проектируется с гранёного кубооктаэдра площади $A_{\mathrm{cuboct}} = (6+2\sqrt{3})a^2$ на вписанную сферу площади $A_{\mathrm{sphere}} = 2\pi a^2$, что умножает $|\square|/N_{\mathrm{today}}^2$ на $\eta_{\mathrm{surf}} = \pi/(3+\sqrt{3})$.

\textbf{Проверка размерностей.}
Проверим размерности каждого звена в системе СИ ($\mathrm{L, M, T}$):
\[
  [\eta_B] = \left[\frac{n_B}{n_\gamma}\right] = \frac{\mathrm{L^{-3}}}{\mathrm{L^{-3}}} = [\alpha_{\mathrm{fs}}^4] = 1,\qquad
  [N_{\mathrm{today}}] = \left[\frac{t_{\mathrm{age}}}{\htick}\right] = \frac{\mathrm{T}}{\mathrm{T}} = 1,
\]
\[
  [\Lambda_{\mathrm{bubble}}] = \left[\frac{\eta_{\mathrm{surf}}|\square|}{N_{\mathrm{today}}^2\,\ell_P^2}\right] = \mathrm{L^{-2}}
  \quad\Longrightarrow\quad
  [\Lambda_{\mathrm{bubble}}\,\ell_P^2] = \mathrm{L^{-2}\cdot L^2} = 1.
\]

\subsubsection*{Ход эксперимента}
Воспроизведение: \texttt{python verify\_principles.py --suite sm} (модуль \texttt{mt\_ca/si\_sm\_rows.py}).

\begin{enumerate}[label=\textbf{Шаг \arabic*:},leftmargin=*,itemsep=0.35em]
\item Подставить upstream-значение $\alpha_{\mathrm{fs}}^{-1} = 137{,}03599919$ в формулу $\eta_B = \frac{45}{208}\alpha_{\mathrm{fs}}^4$.
\item Вычислить $\htick = t_P/\sqrt{2} \approx 3{,}812\times 10^{-44}\ \mathrm{с}$ и $N_{\mathrm{today}} = \lfloor t_{\mathrm{age}}/\htick\rfloor \approx 1{,}141\times 10^{61}$.
\item Вычислить $\Lambda_{\mathrm{bubble}}\ell_P^2 = \eta_{\mathrm{surf}}\cdot 6 / N_{\mathrm{today}}^2$ и рассчитать~$|E_n|$ против Planck~2018.
\end{enumerate}

\begin{table}[htbp]
\centering
\footnotesize
\setlength{\tabcolsep}{4pt}
\caption{Строгая upstream-сверка $\eta_B$, $N_{\mathrm{today}}$ и $\Lambda_{\mathrm{bubble}}\ell_P^2$ с данными Planck 2018 (CE-C-03)}
\label{tab:exp-ce-c-03}
\begin{tabular}{@{}lrrrr@{}}
\toprule
Величина & Формула ДКМ & Модель & Planck 2018 ($1\sigma$) & $|E_n|$ ($k=2$) \\
\midrule
$\eta_B = n_B/n_\gamma$ & $\frac{45}{208}\,\alpha_{\mathrm{fs}}^4$ & $6{,}135\times 10^{-10}$ & $(6{,}12 \pm 0{,}04)\times 10^{-10}$ & $0{,}19$ \\
$N_{\mathrm{today}}$ (тактов $\htick$) & $\lfloor t_{\mathrm{age}}/(\caKappaGeom t_P)\rfloor$ & $1{,}141\times 10^{61}$ & $(1{,}141 \pm 0{,}002)\times 10^{61}$ & $<0{,}05$ \\
$\Lambda_{\mathrm{bubble}}\,\ell_P^2$ & $\frac{3-\sqrt{3}}{2}\frac{2\pi}{N_{\mathrm{today}}^2}$ & $3{,}06\times 10^{-122}$ & $(2{,}89 \pm 0{,}09)\times 10^{-122}$ & $0{,}94$ \\
\bottomrule
\end{tabular}
\end{table}

\subsubsection*{Анализ, расчёт погрешности и выводы}
По закону переноса неопределённости GUM (\cref{eq:gum-relative}) относительная погрешность $\eta_B$ от входа $\alpha_{\mathrm{fs}}$ равна $4\,\delta\alpha/\alpha < 10^{-8}$, что пренебрежимо мало по сравнению с погрешностью Planck~2018 ($0{,}65\%$, тип~B).
Для $\Lambda_{\mathrm{bubble}}\ell_P^2$ относительная погрешность $2\,\delta t_{\mathrm{age}}/t_{\mathrm{age}} \approx 0{,}3\%$ также лежит внутри астрофизической полосы ($\pm 3{,}1\%$), давая $|E_n| = 0{,}94 < 1$.

\textbf{Вывод.}
Включение сферического фактора проектора~$\eta_{\mathrm{surf}}$ полностью замыкает космологический блок ($\eta_B$, $N_{\mathrm{today}}$, $\Lambda_{\mathrm{bubble}}$) внутри $1\sigma$--$2\sigma$ интервала Planck~2018.
\textbf{Статус записи:} опыт закрыт по табл.~\ref{tab:exp-ce-c-03}.

\end{labrecord}

\end{labseries}


\section{Критерии фальсификации}
\label{sec:cosmo-falsification}

Все структурные числа космологического тома однозначно выведены из единственного корня уравнения Каталана $1+2^d=d^2 \iff d=3$ (\texttt{catalan\_\allowbreak lock\_\allowbreak iff\_\allowbreak d3}), поэтому модель не допускает перекалибровки и опровергается в любом из следующих случаев:
\begin{enumerate}[leftmargin=*]
\item \textbf{Нарушение 9-битного порога нуклеации ($\mathrm{IV}(3)=263\in\mathbb{Z}_{512}$):}
  если на трёхмерной решётке~\latGCKshort\ однородный или плоско-стеночный сдвиг фазы без трёхгранного стыка октантов~$K_3$ рождает ненулевой топологический заряд $b=1$, либо если инверсия старшего бита хиральности $b_{\mathrm{chir}}\colon 1\mapsto 0$ в сиде $\mathrm{IV}=263\mapsto 7$ не обращает знак суммарного топологического заряда $\sum n$ (CE-C-01).
\item \textbf{Теорема запрета частиц «тёмной материи» ($D_{\mathrm{dark}}\equiv 0$), отношение ореола $\Omega_{\mathrm{halo}}/\Omega_b = 16/3$ и барионная асимметрия $\eta_B$:}
  по теореме~\ref{thm:cosmo-no-dark-matter-and-halo-ratio} (\texttt{no\_\allowbreak dark\_\allowbreak matter\_\allowbreak and\_\allowbreak vacuum\_\allowbreak holonomy\_\allowbreak halo}) гравитационный импеданс $D(x)$ исчерпывается барионной амплитудой~$\mu$ и калибровочным дефицитом Бома $\delta K_{xy}=\mu_x\mu_y|A_{xy}|$ на единичном полу вакуума $\mu_{\mathrm{vac}}=1$ с точным отношением ёмкости кубооктаэдрального ореола к $3\mathrm{D}$-ядру $2048/384 = 16/3 \approx 5{,}333$. Модель даёт строгий нулевой прогноз ($N_{\mathrm{DM\text{-}particles}}\equiv 0$) для всех детекторов частиц тёмной материи и опровергается, если любой эксперимент обнаружит калибровочно-стерильную частицу тёмной материи ($\mu>1$, $\delta K_{xy}=0$) либо если космологические измерения выведут $\eta_B$ за пределы $\frac{45}{208}\alpha_{\mathrm{fs}}^4 \approx 6{,}135\times 10^{-10}$ или отношение $\Omega_{\mathrm{halo}}/\Omega_b$ за пределы $16/3$ (CE-C-02, CE-C-03).
\item \textbf{Горизонтная кривизна $\Lambda_{\mathrm{bubble}}$:}
  если космологический член окажется независимым от каузального возраста пузыря $\Lambda_{\mathrm{bubble}}\ell_P^2 = \eta_{\mathrm{surf}}|\square|/N_{\mathrm{today}}^2 \approx 3{,}06\times 10^{-122}$ (CE-C-03).
\item \textbf{Естественный тепловой квант реликта $\Theta_{\mathrm{CMB}}$ и двухкаскадное горлышко $\tau_{\mathrm{neck}} = \phaseRingCount/2 = 256$:}
  если отношение энергии ионизации водорода к тепловому кванту рекомбинации отклонится от $E_{\mathrm{Ry}}/\Theta_{\mathrm{rec}} = \frac{5}{2}\ln(208/(45\alpha_{\mathrm{fs}}^4)) \approx 53{,}030$, либо если при прямом 3D-прогоне сида $\mathrm{IV}=263$ разрушится полупериодное горлышко $\tau_{\mathrm{neck}}=256$ с рождением трёхквантового ядра $|w|=3$ на соседнем ребре~$N_{12}$ выжившего реликтового зерна (CE-C-04).
\end{enumerate}
